\documentclass[11pt,a4paper]{article}
\usepackage[margin=0.85in,top=1in,bottom=1in]{geometry}
\usepackage{amsmath,amssymb,amsfonts}
\usepackage{graphicx}
\usepackage{float}
\usepackage{enumitem}
\usepackage{microtype}
\usepackage{caption}
\usepackage[hidelinks]{hyperref}
\usepackage[most,breakable]{tcolorbox}
\tcbuselibrary{breakable,skins}
\usepackage{fontspec}
\setmainfont{DejaVu Serif}
\setsansfont{DejaVu Sans}
\setmonofont{DejaVu Sans Mono}
\captionsetup{font=small,labelfont=bf,skip=4pt}
\setlist{leftmargin=1.5em,topsep=2pt}

%% Breakable problem card — prevents content cutoff across pages
\newtcolorbox{solcard}[3]{
  breakable, enhanced,
  colback=white, colframe=gray!50, arc=2pt,
  title={\textbf{#1}\hfill\textsf{\small #3\,pts}},
  fonttitle=\normalsize\sffamily\bfseries,
  before upper={\small\textit{#2}\par\smallskip\hrule height 0.4pt\smallskip},
  top=4pt, bottom=6pt, left=7pt, right=7pt,
  parbox=false,
}

\title{\textbf{Spaceship in the atmosphere --- Solution}}
\author{\normalsize X25 (2025) — English translation}
\date{}
\begin{document}\maketitle\vspace{-0.5em}
\begin{solcard}{A1}{First, we obtain the exact equations that describe the dynamics of the spacecraft. Express the time derivatives of altitude $h$ and range $s$ in terms of $v$, $\gamma$, $r_0$, $r$.}{0.20}

The rate of decrease in height is the vertical projection of speed $v \sin \gamma$, the speed of horizontal movement $v \cos \gamma$ after projection onto the surface of the Earth decreases by $r/r_0$ times.

\par\smallskip\noindent\textbf{Answer:} \textbackslash{}textbf{Answer:} $$\dot h=-v \sin{\gamma}, \quad\dot s= v\cos{\gamma}\cdot r_0/r$$
\end{solcard}

\begin{solcard}{A2}{Find the derivatives of the velocity modulus $v$ and angle $\gamma$. The answer may also include the lift force $L$, the air resistance force $D$, the acceleration due to gravity at a given height $g$, the mass of the spacecraft $m$, $v$, $\gamma$, $r$.}{1.00}

Newton's second law in projection onto the axes along and perpendicular to the speed: $$ma_\tau=-D+mg\sin{\gamma}, \quad ma_n=-L+mg\cos{\gamma}.$$Tangential acceleration  $a_\tau=\dot v$, therefore $$ a_\tau = - \frac{D}{m} + g \sin\gamma. $$Normal acceleration $a_n = \omega v $, where $\omega$ – angular velocity поворота vectorа скорости, которую считаем положительной, if angular velocity поворачивается внfrom.  Эта velocity не equals производной угла $\gamma$, since направление горизонтали меняется от точки к точке с угловой velocityю $\omega_h = v \cos \gamma/r$, then $$ \omega = \dot{\gamma} + \frac{v \cos \gamma}{r}, $$therefore equation for нормального ускорения $$ mv \left (\dot{\gamma} + \frac{v \cos \gamma}{r}\right) = - L + mg \cos \gamma, $$where do we find the required derivative

\par\smallskip\noindent\textbf{Answer:} \textbackslash{}textbf{Answer:} $$\dot v= -\frac{D}{m}+g\sin{\gamma}, \quad \dot\gamma=-\frac{L}{mv}+\frac{g}{v}\cos{\gamma}-\frac{v\cos{\gamma}}{r}$$
\end{solcard}

\begin{solcard}{B1}{Obtain an expression for the derivative $dv/dh$. Represent the answer as $$ \frac{dv}{dh} = A f(v,\,h), $$ where $A$ – постоянная, которая может зависеть от всех постоянных, характеризующих космический корабль and атмосферу ($m$, $S$, $C_D$, $C_L$, $\rho_0$, $\beta$, $\gamma$), а $f$ – некоторая function высоты and скорости. В дальнейшем используйте постоянную $\mathbf{A}$ to write answers.}{0.40}

Let us substitute the expression for the air resistance force $D$ into the equation of motion and obtain, neglecting gravitational acceleration $$\dot v = - \dfrac{\rho_0 SC_D}{2m}v^2e^{-\beta h}.$$Taking into account формулы for $\dot{h} = - v\sin \gamma$ we get $$\dfrac{dv}{dh}=\dfrac{\dot v}{\dot h}=\dfrac{\rho_0 C_D S}{2m\sin{\gamma}}\cdot ve^{-\beta h}$$

\par\smallskip\noindent\textbf{Answer:} \textbackslash{}textbf{Answer:} $$\dfrac{dv}{dh}=A ve^{-\beta h}, \quad A = \dfrac{\rho_0 C_D S}{2m\sin{\gamma}}$$
\end{solcard}

\begin{solcard}{B2}{Find the dependence of the speed of the spacecraft on altitude. At the initial height $h_0$ the speed was equal to $v_0$. Express your answer in terms of $A$, $\beta$, $h$, $v_0$, $h_0$. Also get an approximate answer, assuming that at an altitude of $h_0$ the atmospheric density is negligible (in this case the answer should not include $h_0$). In the following, use an approximate expression throughout.}{0.80}

Dividing the variables, we get $$ \frac{dv}{v} = A e^{- \beta h} dh, $$проинтегрировав, we get $$\int\limits_{v_0}^{v}\dfrac{dv}{v}=\int\limits_{h_0}^{h}Ae^{-\beta h}dh, \quad \ln\dfrac{v}{v_0}=-\dfrac{A}{\beta}(e^{-\beta h}-e^{-\beta h_0}).$$Hence let us find expression for скорости $$v=v_0\exp\left(-\dfrac{A}{\beta}(e^{-\beta h}-e^{-\beta h_0})\right)$$For/At достаточно большом значении $h_0$ terms $e^{- \beta h_0} \ll e^{- \beta h}$ можно пренебречь $$v=v_0\exp\left(-\dfrac{A}{\beta}e^{-\beta h}\right)$$

\par\smallskip\noindent\textbf{Answer:} \textbackslash{}textbf{Answer:} $$v=v_0\exp\left(-\dfrac{A}{\beta}(e^{-\beta h}-e^{-\beta h_0})\right), \quad v\approx v_0\exp\left(-\dfrac{A}{\beta}e^{-\beta h}\right)$$
\end{solcard}

\begin{solcard}{B3}{Find the dependence of the acceleration $a$ (that is, the derivative of the velocity modulus) of the spacecraft on altitude. Express your answer in terms of $v_0$, $A$, $\beta$, $h$, $\gamma$.}{0.30}

Let us substitute the formula for velocity into the equation for tangential acceleration and obtain $$a=\dot v=-\frac{D}{m}=- \dfrac{\rho_0 SC_D}{2m}v^2e^{-\beta h}=-Av_0^2\sin{\gamma}\exp\left(-\dfrac{2A}{\beta}e^{-\beta h}-\beta h\right)$$

\par\smallskip\noindent\textbf{Answer:} \textbackslash{}textbf{Answer:} $$a=-Av_0^2\sin{\gamma}\exp\left(-\dfrac{2A}{\beta}e^{-\beta h}-\beta h\right)$$
\end{solcard}

\begin{solcard}{B4}{Find the height $h_c$ at which the acceleration modulus is maximum. Obtain a formula for the maximum value of acceleration modulus $a_{\text{max}}$ and velocity $v_с$ at height $h_с$. Express your answers using $A$, $\beta$, $\gamma$, $v_0$.}{0.50}

The acceleration module is maximum, with a minimum exponent, $$\dfrac{d}{dh}\left(-\dfrac{2A}{\beta}e^{-\beta h}-\beta h\right)=2A e^{- \beta h} - \beta=0.$$Hence критическое value высоты $$h_c=\dfrac{1}{\beta}\ln\dfrac{2A}{\beta}.$$Substituting это value высоты в формулы for высоты and ускорения, we get $$a_{max}=Av_0^2\sin{\gamma}\exp\left(-\dfrac{2A}{\beta}e^{-\beta h_c}-\beta h_c\right)=\dfrac{\beta v_0^2\sin{\gamma}}{2e}$$$$v_c=v_0\exp\left(-\dfrac{A}{\beta}e^{-\beta h_c}\right)=v_0e^{-1/2}$$

\par\smallskip\noindent\textbf{Answer:} \textbackslash{}textbf{Answer:} $$h_c=\dfrac{1}{\beta}\ln\dfrac{2A}{\beta}, \quad v_c=v_0e^{-1/2}, \quad a_{max}=\dfrac{\beta v_0^2\sin{\gamma}}{2e}$$
\end{solcard}

\begin{solcard}{B5}{For the given numerical data, find the speed (in km/s) and acceleration (in units of $g_0$) at heights $h_1 = 80~\text{km}$, $h_2 = 60~\text{km}$, $h_3 = 40~\text{km}$.}{0.60}

Substituting the height values ​​into the formulas for speed and acceleration, we get

\par\smallskip\noindent\textbf{Answer:} \textbackslash{}textbf{Answer:} $h, ~\mathrm{km}$ $a/g_0$ $v, ~\mathrm{km}/с$ $80$ $-0.20$ $7.7$ $60$ $-2.8$ $7.2$ $40$ $-5.9$ $2.7$
\end{solcard}

\begin{solcard}{B6}{Find the maximum acceleration (in units of $g_0$) and the critical height $h_c$ at which it is achieved.}{0.20}
\par\smallskip\noindent\textbf{Answer:} \textbackslash{}textbf{Answer:} $$a_{max}=8.6~g_0, \quad h_c=45.5~\mathrm{km}$$
\end{solcard}

\begin{solcard}{C1}{Find the first escape velocity $v_s$ – the speed of the spacecraft in a circular orbit, the radius of which is equal to the radius of the Earth. Express your answer in terms of $r_0$, $g_0$.}{0.10}

Let us write the equation of motion for acceleration when moving in a circle $$ \frac{v^2}{r_0} = g_0, $$

\par\smallskip\noindent\textbf{Answer:} \textbackslash{}textbf{Answer:} $$v_s=\sqrt{g_0r_0}$$
\end{solcard}

\begin{solcard}{C2}{Determine the speed at which the spacecraft must move at altitude $h$ for the motion described in the introduction to this part. Express your answer in terms of $v_s$, $\rho_0$, $\beta$, $C_L$, $S$, $m$, $r_0$, $h$.}{0.40}

Let us write the equation for normal acceleration, taking into account that the derivative $\dot{\gamma}$ is small, and $\cos \gamma \approx 1$ and it can be neglected $$-\frac{L}{m}+g_0 -\frac{v^2}{r_0} = 0.$$Substituting формулу for подъемной силы, а also for зависимости плотности воздуха от высоты, we get $$ v^2 \left( \frac{1}{r_0} + \frac{C_L S \rho_0}{2 m}e^{- \beta h}\right) = g_0, $$hence let us find $$v^2=\frac{g_0r_0}{1+\dfrac{C_L \rho_0 r_0 S}{2m}\cdot e^{-\beta h}}$$

\par\smallskip\noindent\textbf{Answer:} \textbackslash{}textbf{Answer:} $$v= \frac{v_s}{\sqrt{1+\dfrac{C_L \rho_0 r_0 S}{2m}\cdot e^{-\beta h}} }$$
\end{solcard}

\begin{solcard}{C3}{Determine the acceleration $a$ (derivative of the velocity modulus) created by the force of air resistance during such movement. Express your answer in terms of $v$, $C_L$, $C_D$, $r_0$, $g_0$.}{0.40}

Lift and air resistance are related by the relation $$ \frac{D}{L} = \frac{C_L}{C_D} $$ $$\dot v=-D/m=-\dfrac{C_D}{C_Lm}\cdot L=-\dfrac{C_D}{C_L}\left(g_0-\dfrac{v^2}{r_0}\right)$$

\par\smallskip\noindent\textbf{Answer:} \textbackslash{}textbf{Answer:} $$a=-\dfrac{C_D}{C_L}\left(g_0-\dfrac{v^2}{r_0}\right)$$
\end{solcard}

\begin{solcard}{C4}{Let the initial speed of the spacecraft be $v_1< v_s$ and the final speed $v_2 < v_1$. Find the horizontal displacement of the ship $s$ during the motion. Consider that the entire change in the velocity module occurs due to the force of air resistance. Express your answer using $r_0$, $C_L$, $C_D$, $v_s$, $v_1$, $v_2$.}{0.50}

Using the relations $a = dv/dt$ and $v = ds/dt$ (in the approximation $\cos \gamma \approx 1$), we obtain the connection between displacement and change in speed, $$ \frac{ds}{dv} = \frac{v}{a} = - \frac{C_L}{C_D} \frac{v}{g_0 - v^2/r_0} =  - \frac{C_L}{C_D} r_0 \frac{v}{v_s^2 - v^2}. $$Integrating, we get $$ s = - \frac{C_L}{C_D} r_0 \int\limits_{v_1}^{v_2} \frac{vdv}{v_s^2 -v^2} = - \frac{C_L}{2C_D} r_0  \int\limits_{v_1}^{v_2} dv\left( \frac{1}{v_s - v} + \frac{1}{v_s + v}\right) = \frac{C_L}{2 C_D}r_0 \ln \frac{v_s^2 -v_2^2}{v_s^2 -v_1^2}. $$

\par\smallskip\noindent\textbf{Answer:} \textbackslash{}textbf{Answer:} $$s=\dfrac{C_L}{2C_D}r_0\ln\dfrac{v_s^2-v_2^2}{v_s^2-v_1^2}$$
\end{solcard}

\begin{solcard}{C5}{Let the following parameters be given for the spaceship: $m  = 84 \times 10^3~\text{kg}$ $S = 250~\text{m}^2$ $C_D = 0.8$ $C_L = 0.9$ Find the values ​​of speed (in km/s) and acceleration due to air resistance (in units of $g_0$) at altitudes $h_1 = 80~\text{km}$, $h_1 = 60~\text{km}$, $h_1 = 40~\text{km}$.}{0.60}

Find the values ​​of speed (in km/s) and acceleration due to air resistance (in units of $g_0$) at heights $h_1 = 80~\text{km}$, $h_1 = 60~\text{km}$, $h_1 = 40~\text{km}$.


Using formulas for speed and acceleration expressed in terms of speed, we obtain

\par\smallskip\noindent\textbf{Answer:} \textbackslash{}textbf{Answer:} $h, ~\mathrm{km}$ $a/g_0$ $v, ~\mathrm{km}/c$ $80$ $-0.13$ $7.3$ $60$ $-0.65$ $4.1$ $40$ $-0.87$ $1.2$
\end{solcard}

\begin{solcard}{С6}{Under the conditions of the previous paragraph, calculate the range of horizontal movement of the spacecraft when lowering from height $h_0 = 90~\text{km}$ to height $h_f = 30~\text{km}$.}{0.20}

Let's calculate the initial and final speeds of the spacecraft $$v_1 = v(h_0)=7.7~ км/с, \quad v_2 = v(h_f)=0.61~\mathrm{km}/c,$$then the range of movement

\par\smallskip\noindent\textbf{Answer:} \textbackslash{}textbf{Answer:} $$s=11.4\cdot 10^3~ км$$
\end{solcard}

\begin{solcard}{D1}{Write down expressions for energy $E$ and angular momentum $L$. Express your answer using $G, m, a, e, M$. Note: It may be convenient to write the ZSE at the apocentric or periapsis point.}{0.40}

Note: It may be convenient to write the ZSE at the apocentric or periapsis point.


Let the distances from the ground to the ship at the pericenter (the closest point) and apocenter (the most distant point) be equal to $r_1$ and $r_2$, respectively. The speeds of the ship at these points are $L/(mr_1)$ and $L/(mr_2)$, respectively. Then: $$E =  \frac{L^2}{2 mr_1^2} - \frac{GMm}{r_1} =   \frac{L^2}{2 m r_2^2} - \frac{GMm}{r_2},$$Let us express hence angular momentum  $L$: $$ L^2 = 2 m r_1^2 E_1+ 2 GMm^2 r_1 = 2 m r_2^2 E_2 + 2 GMm^2 r_2. $$Hence let us express энергию $$ E (r_2^2 - r_1^2) = GMm^2 (r_1  - r_2), $$and therefore $$ E = - \frac{GMm}{r_1 + r_2} = - \frac{GMm}{2a}. $$Таким же обсом let us express angular momentum $$ \frac{L^2}{2 m} \left( \frac{1}{r_1^2} - \frac{1}{r_2^2}\right) = GMm \left(\frac{1}{r_1} - \frac{1}{r_2} \right), \quad L^2 = 2G Mm^2 \frac{r_1 r_2}{r_1 + r_2}. $$Expressing расстояния via/through большую полуось and эксцентриситет $r_1 = a (1-e)$, $r_2 = a (1+ e)$ let us find $$ L^2 = G Mm^2 a (1-e^2). $$

\par\smallskip\noindent\textbf{Answer:} \textbackslash{}textbf{Answer:} $$E=-\dfrac{GMm}{2a}, \quad L = \sqrt{GMm^2a(1-e^2)}$$
\end{solcard}

\begin{solcard}{D2}{Write down the moment equation and expression for the power of the friction force $P$. Express your answer using $\vec r, \vec v, \alpha$.}{0.20}
\par\smallskip\noindent\textbf{Answer:} \textbackslash{}textbf{Answer:} $$\frac{d}{dt} \vec{L}= -\alpha [\vec r, \vec v]$$ $$P=-\alpha v^2$$
\end{solcard}

\begin{solcard}{D3}{Obtain an expression for the dependence of $\vec L$ on $t$ if at time $t=0$ $~\vec L=\vec L_0$. Express your answer using $\alpha$, $m$, $\vec L_0$.}{0.20}

Let's use the definition for $\vec L$ and write the moment equation in the form $$\dfrac{d}{dt}\vec L=-\dfrac{\alpha}{m}\vec L,$$

\par\smallskip\noindent\textbf{Answer:} \textbackslash{}textbf{Answer:} $$\vec L=\vec L_0 \exp\left(-\dfrac{\alpha}{m}t\right)$$
\end{solcard}

\begin{solcard}{D4}{Show that: $$[\vec a, \vec L]=\beta \frac{d}{dt}\vec e_i$$ Where $\vec a$ – acceleration ракеты, $\vec e_i$ – некоторый единичный vector полярных координат с началом в центре Земли. Find $\vec e_i$ and $\beta$. For/Atмечание: Может быть удобным использовать expression $\vec L=mr^2 \vec\omega$.}{0.30}

Note: It may be convenient to use the expression $\vec L=mr^2 \vec\omega$.


Newton's second law is $$m\vec a=-\dfrac{GmM}{r^2}\vec e_r.$$Substituting acceleration в vectorное произведение: $$[\vec a, ~\vec L]=-\left[\dfrac{GM}{r^2}\vec e_r, ~\vec L\right]= -GM\left[\dfrac{1}{r^2}\vec e_r, ~mr^2\vec\omega\right] =GmM\left[\vec\omega,~\vec e_r\right]=GmM\dfrac{d}{dt}\vec e_r.$$

\par\smallskip\noindent\textbf{Answer:} \textbackslash{}textbf{Answer:} $$\beta=GMm, \quad \vec e_i=\vec e_r$$
\end{solcard}

\begin{solcard}{D5}{Multiplying expression (1) by $\vec r$ scalarly and using the results of point D4, obtain the dependence of the modulus of the radius vector $\vec r$ on $\varphi$ – the angle between vectors $\vec A$ and $\vec e_r$. Express your answer using $L, m, G, e, \varphi, M$.}{0.30}

Multiply the expression by $\vec{r}$, we get $$ (\vec r, [\vec v, \vec L]) = GMmr + \vec{A} \cdot \vec{r}. $$В левой части переставим множители, using свойства смешанного произведения: $$ (\vec r, [\vec v, \vec L]) = (\vec{L}, [\vec{r}, \vec{v}]) = \frac{L^2}{m}. $$Then we get $$ \frac{L^2}{m} = GmMr(1 + e \cos \varphi), $$where we find

\par\smallskip\noindent\textbf{Answer:} \textbackslash{}textbf{Answer:} $$r=\dfrac{L^2/GMm^2}{1+e\cos{\varphi}}$$
\end{solcard}

\begin{solcard}{D6}{Using expression (1) and the results of point D4, obtain the dependence of the square of the speed $v^2$ on $\varphi$. Express your answer using $L, G, m, e, \varphi, M$.}{0.50}

Let's square expression (1): $$[\vec v, ~\vec L]^2=v^2L^2= (GmM)^2 + 2 GmM \vec{e}_r \cdot \vec{A} + A^2 = (GmM)^2 (1 + 2e \cos \varphi + e^2),$$

\par\smallskip\noindent\textbf{Answer:} \textbackslash{}textbf{Answer:} $$v^2=\dfrac{(GMm)^2}{L^2}(1+e^2+2e\cos{\varphi})$$
\end{solcard}

\begin{solcard}{D7}{Obtain an expression for the time-average squared velocity $\langle v^2\rangle $ as an integral over the angle $\varphi$. Express your answer in terms of $  L,  m,  e,  M, G, a, \varphi, \tau$, the rocket’s orbital period. Note: It may be convenient to express the square of the average velocity from the averaged ZSE ($\langle T+U\rangle=\langle E\rangle$, where $T$ is kinetic energy, $U$ is potential).}{0.40}

Note: It may be convenient to express the square of the average velocity from the averaged ZSE ($\langle T+U\rangle=\langle E\rangle$, where $T$ is kinetic energy, $U$ is potential energy).


Method 1: By definition, the mean square of velocity is expressed through the integral as ($\tau$ – period of motion) $$\langle v^2\rangle= \dfrac{1}{\tau}\int\limits_{0}^{\tau}v^2 dt.$$Дифференциалы времени and угла $\varphi$ можно связать с помощью momentа momentumа: $$L=mr^2\dfrac{d\varphi}{dt}.$$Перейдем к интегрированию по углу $$\langle v^2 \rangle =\dfrac{1}{\tau}\int\limits_{0}^{\tau}v^2\dfrac{dt}{d\varphi}d\varphi=\dfrac{m}{\tau L}\int\limits_{0}^{2\pi}v^2r^2d\varphi.$$Substituting выражения from D5, D6, we get $$\langle v^2 \rangle= \dfrac{L}{m\tau} \int\limits_{0}^{2\pi}\dfrac{1+e^2+2e\cos{\varphi}}{(1+e\cos{\varphi})^2}d\varphi .$$. Метод 2: From усредненного ЗСЭ we get: $$\langle T\rangle=\dfrac{m}{2}\langle v^2\rangle=\langle E\rangle-\langle U\rangle=E+\left\langle \dfrac{GMm}{r}\right\rangle$$Let us write expression for среднего по времени $1/r$ в виде интеграла по углу: $$\left\langle \dfrac{1}{r}\right\rangle= \dfrac{1}{\tau}\int\limits_{0}^{\tau}\dfrac{1}{r} dt = \dfrac{1}{\tau}\int\limits_{0}^{\tau}\dfrac{1}{r} dt\dfrac{dt}{d\varphi}d\varphi=\dfrac{m}{\tau L}\int\limits_{0}^{2\pi}rd\varphi $$Substituting все в expression for $\langle v^2 \rangle$: $$\langle v^2 \rangle=\dfrac{2}{m}\left(E+\dfrac{L}{\tau}\int\limits_{0}^{2\pi}\dfrac{d\varphi}{1+e\cos{\varphi}}\right)$$

\par\smallskip\noindent\textbf{Answer:} \textbackslash{}textbf{Answer:} $$\langle v^2 \rangle= \dfrac{L}{m\tau} \int\limits_{0}^{2\pi}\dfrac{1+e^2+2e\cos{\varphi}}{(1+e\cos{\varphi})^2}d\varphi$$.
\end{solcard}

\begin{solcard}{D8}{Obtain an expression for $ \langle\dot E\rangle $. Express your answer using $G, m, a, e, M$.}{1.00}

Method 1: $$\langle v^2 \rangle= \dfrac{L}{m\tau} \int\limits_{0}^{2\pi}\dfrac{1+e^2+2e\cos{\varphi}}{(1+e\cos{\varphi})^2}d\varphi$$. Let's use the universal trigonometric substitution: $$\int\limits_{0}^{2\pi}\dfrac{1+e^2+2e\cos{\varphi}}{(1+e\cos{\varphi})^2}d\varphi=2\int\limits_{0}^{+\infty}\dfrac{1+e^2+2e(1-t^2)/(1+t^2)}{(1+e(1-t^2)/(1+t^2))^2}\dfrac{2dt}{1+t^2}=$$ $$=4\int\limits_{0}^{+\infty}\dfrac{(1+e^2)(1+t^2)+2e(1-t^2)}{(1+t^2+e(1-t^2))^2}dt=4\int\limits_{0}^{+\infty}\dfrac{(1+e)^2+t^2(1-e)^2}{(1+e+t^2(1-e))^2}dt=$$ $$=4\int\limits_{0}^{+\infty}\dfrac{(1+e)^2+t^2(1-e)^2}{(1+e+t^2(1-e))^2}dt=4\int\limits_{0}^{+\infty}\dfrac{1+t^2(1-e)^2/(1+e)^2}{(1+t^2(1-e)/(1+e))^2}dt$$ Перейдем к переменной $\tan u=t\sqrt{\dfrac{1-e}{1+e}}\Rightarrow \dfrac{du}{\cos^2u}=dt\sqrt{\dfrac{1-e}{1+e}}$ $$4\int\limits_{0}^{+\infty}\dfrac{1+t^2(1-e)^2/(1+e)^2}{(1+t^2(1-e)/(1+e))^2}dt=4\sqrt{\dfrac{1+e}{1-e}}\int\limits_{0}^{+\pi/2}\dfrac{1+\tan^2u(1-e)/(1+e)}{(1+\tan^2u)^2}\dfrac{du}{\cos^2u}=$$ $ $=\textbackslash{}dfrac{4}{\textbackslash{}sqrt{1-e^2}}\textbackslash{}int\textbackslash{}limits_{0}^{+\textbackslash{}pi/2}((1+e)\textbackslash{}cos^2u+\textbackslash{}sin^2u(1-e))du=\textbackslash{}dfrac{4}{\textbackslash{}sqrt{1-e^2}}\textbackslash{}int\textbackslash{}limits_{0}^{+\textbackslash{}pi/2}du(1+e\textbackslash{}cos2u)=\textbackslash{}dfrac{2\textbackslash{}pi}{\textbackslash{}sqrt{1-e^2}}$$ Подставим в выражение для среднего: $$\textbackslash{}langle v^2 \textbackslash{}rangle= \textbackslash{}dfrac{L}{m\textbackslash{}tau} \textbackslash{}dfrac{2\textbackslash{}pi}{\textbackslash{}sqrt{1-e^2}}=\textbackslash{}sqrt{GMa}\textbackslash{}cdot\textbackslash{}sqrt{\textbackslash{}dfrac{GM}{a^3}}=\textbackslash{}dfrac{GM}{a}$$ Метод 2: $$\textbackslash{}langle v^2 \textbackslash{}rangle=\textbackslash{}int\textbackslash{}limits_{0}^{\textbackslash{}tau}\textbackslash{}dfrac{(GMm)^2}{L^2\textbackslash{}tau}(1+e^2+2e\textbackslash{}cos{\textbackslash{}varphi})dt$$ Усреднение константы даст саму константу, а из результата предыдущего пункта ясно, что переход к переменной $\varphi$ даст множитель $1/(1+e\cos\varphi)^2 $, соберем в числителе $(1+e\cos\varphi)$: $$=\textbackslash{}langle v^2 \textbackslash{}rangle=\textbackslash{}int\textbackslash{}limits_{0}^{\textbackslash{}tau}\textbackslash{}dfrac{(GMm)^2}{L^2\textbackslash{}tau}(e^2-1+2(1+e\textbackslash{}cos{\textbackslash{}varphi}))dt=$$ $$=\textbackslash{}dfrac{(GMm)^2}{L^2}(e^2-1)+ \textbackslash{}int\textbackslash{}limits_{0}^{\textbackslash{}tau}\textbackslash{}dfrac{(GMm)^2}{L^2\textbackslash{}tau}(2(1+e\textbackslash{}cos{\textbackslash{}varphi}))dt=$$ $$=\textbackslash{}dfrac{(GMm)^2}{L^2}(e^2-1)+\textbackslash{}int\textbackslash{}limits_{0}^{\textbackslash{}tau}\textbackslash{}dfrac{(GMm)^2}{L^2\textbackslash{}tau}(1+e\textbackslash{}cos{\textbackslash{}varphi})dt=$$ Возьмем получившийся интеграл при помощи универсальной тригонометрической подстановки: $$\textbackslash{}dfrac{(GMm)^2}{L^2\textbackslash{}tau}\textbackslash{}int\textbackslash{}limits_{0}^{\textbackslash{}tau}(1+e\textbackslash{}cos{\textbackslash{}varphi})dt=\textbackslash{}dfrac{2L}{m\textbackslash{}tau}\textbackslash{}int\textbackslash{}limits_{0}^{2\textbackslash{}pi}\textbackslash{}dfrac{d\textbackslash{}varphi}{1+e\textbackslash{}cos{\textbackslash{}varphi}}=\textbackslash{}dfrac{4L}{m\textbackslash{}tau}\textbackslash{}int\textbackslash{}limits_{0}^{+\textbackslash{}infty}\textbackslash{}dfrac{1}{1+e(1-t^2)/(1+t^2)}\textbackslash{}dfrac{2dt}{1+t^2}=$⟦2 5⟧$=\textbackslash{}dfrac{8L}{m\textbackslash{}tau}\textbackslash{}int\textbackslash{}limits_{0}^{+\textbackslash{}infty}\textbackslash{}dfrac{1}{1+t^2+e(1-t^2)}dt=\textbackslash{}dfrac{8L}{m\textbackslash{}tau}\textbackslash{}int\textbackslash{}limits_{0}^{+\textbackslash{}infty}\textbackslash{}dfrac{dt}{1+e+t^2(1-e)}=\textbackslash{}dfrac{8L}{m\textbackslash{}tau}\textbackslash{}dfrac{\textbackslash{}pi/2}{\textbackslash{}sqrt{1-e^2}}$$ Подставим в выражение для средней скорости: $$\textbackslash{}langle v^2 \textbackslash{}rangle= \textbackslash{}dfrac{(GMm)^2}{L^2}(e^2-1)+\textbackslash{}dfrac{4\textbackslash{}pi}{m\textbackslash{}tau}\textbackslash{}dfrac{L}{\textbackslash{}sqrt{1-e^2}}=-\textbackslash{}dfrac{GM}{a}+2\textbackslash{}sqrt{GMa}\textbackslash{}cdot \textbackslash{}sqrt{\textbackslash{}dfrac{GM}{a^3}}=-\textbackslash{}dfrac{GM}{a}+2\textbackslash{}dfrac{GM}{a}=\textbackslash{}dfrac{GM}{a}$$ Метод 3: $$\textbackslash{}langle v^2 \textbackslash{}rangle=\textbackslash{}dfrac{2}{m}\textbackslash{}left(E+\textbackslash{}dfrac{L}{\textbackslash{}tau}\textbackslash{}int\textbackslash{}limits_{0}^{2\textbackslash{}pi}\textbackslash{}dfrac{d\textbackslash{}varphi}{1+e\textbackslash{}cos{\textbackslash{}varphi}}\textbackslash{}right)$$ Возьмем получившийся интеграл при помощи универсальной тригонометрической подстановки: $$\textbackslash{}dfrac{2L}{m\textbackslash{}tau}\textbackslash{}int\textbackslash{}limits_{0}^{2\textbackslash{}pi}\textbackslash{}dfrac{d\textbackslash{}varphi}{1+e\textbackslash{}cos{\textbackslash{}varphi}}=\textbackslash{}dfrac{4L}{m\textbackslash{}tau}\textbackslash{}int\textbackslash{}limits_{0}^{+\textbackslash{}infty}\textbackslash{}dfrac{1}{1+e(1-t^2)/(1+t^2)}\textbackslash{}dfrac{2dt}{1+t^2}=$$ $$=\textbackslash{}dfrac{8L}{m\textbackslash{}tau}\textbackslash{}int\textbackslash{}limits_{0}^{+\textbackslash{}infty}\textbackslash{}dfrac{1}{1+t^2+e(1-t^2)}dt=\textbackslash{}dfrac{8L}{m\textbackslash{}tau}\textbackslash{}int\textbackslash{}limits_{0}^{+\textbackslash{}infty}\textbackslash{}dfrac{dt}{1+e+t^2(1-e)}=\textbackslash{}dfrac{8L}{m\textbackslash{}tau}\textbackslash{}dfrac{\textbackslash{}pi/2}{\textbackslash{}sqrt{1-e^2}}$$ Подставим в выражение для средней скорости: $$\textbackslash{}langle v^2 \textbackslash{}rangle= \textbackslash{}dfrac{2E}{m}+\textbackslash{}dfrac{4\textbackslash{}pi}{m\textbackslash{}tau}\textbackslash{}dfrac{L}{\textbackslash{}sqrt{1-e^2}}=-\textbackslash{}dfrac{GM}{a}+2\textbackslash{}sqrt{GMa}\textbackslash{}cdot \textbackslash{}sqrt{\textbackslash{}dfrac{GM}{a^3}}=-\textbackslash{}dfrac{GM}{a}+2\textbackslash{}dfrac{GM}{a}=\textbackslash{}dfrac{GM}{a}$$ Окончательно: $$\textbackslash{}langle \textbackslash{}dot E\textbackslash{}rangle=-\textbackslash{}dfrac{\textbackslash{}alpha GM}{a}$$


$$\langle \dot E\rangle=-\dfrac{\alpha GM}{a}$$

\par\smallskip\noindent\textbf{Answer:} \textbackslash{}textbf{Answer:} $$\langle \dot E\rangle=-\dfrac{\alpha GM}{a}$$
\end{solcard}

\begin{solcard}{D9}{Using the results of points D1, D3, D8, obtain the dependences of $a$ and $e$ on $t$, if at time $t=0$ $a=a_0$, $e=e_0$.}{0.40}

From the result of the previous point: $$\dot E=-\dfrac{\alpha GM}{a} \Rightarrow \dfrac{d}{dt}\left(\dfrac{1}{a}\right)=\dfrac{2\alpha}{m}\dfrac{1}{a}$$ $$a=a_0e^{-2\alpha t/m}$$ Substituting expression for полуоси в зависимость momentа momentumа от времени: $$L=m\sqrt{GMa(1-e^2)}\Rightarrow \sqrt{a(1-e^2)}=\sqrt{a_0(1-e_0^2)}e^{-\alpha t/m}$$ $$\sqrt{1-e^2}=\sqrt{1-e_0^2}\Rightarrow e=e_0$$

\par\smallskip\noindent\textbf{Answer:} \textbackslash{}textbf{Answer:} $$a=a_0e^{-2\alpha t/m}$$ $$e=e_0$$
\end{solcard}

\begin{solcard}{D10}{Find the time of descent of orbit $T$ from height $h_0=408~\mathrm{km}$ to height $h=400~\mathrm{km}$, if parameter $\alpha=7.17\cdot 10^{-5}~kg/с$, mass $m=420~т$, $r_0=6378~\mathrm{km}$. Express your answer in days.}{0.10}

From the result of the previous point: $$T=\dfrac{m}{2\alpha}\ln\dfrac{r_0+h_0}{r_0+h}$$

\par\smallskip\noindent\textbf{Answer:} \textbackslash{}textbf{Answer:} $$T=40~ суток$$
\end{solcard}

\end{document}